\scanpage{94}
\mspage{23}
de la face. On notera que pour décrire le groupe $\Pi_P^\tau$, et l'équivalence (101), on n'a pas besoin d'autre donnée que du pol. combinatoire $P$, le groupe $\Pi_P^\tau$ est défini simplement par (62) (sans référence au choix d'un point-base $P_\omega$ dans un espace $\Sigma_P^*$ !). Par contre, pour définir $\Pi_P$ et l'équivalence (102), on a besoin du choix d'une orientation $\omega$ de $P$, ce qui permet d'introduire $\Pi_P$ p. ex. comme $\Pi_1(\Sigma_P^*,P_\omega)$. À vrai dire, on pourrait définir $\Pi_P$, sans choix de $\omega$, comme noyau de l'hom. can. $\Pi_P^\tau\unsure{\mathbin{-}|\longrightarrow}\{\pm1\}$, de façon à avoir la suite exacte

\diagram{\begin{tikzcd}[ampersand replacement=\&]
1\arrow[r]\&\Pi_P\arrow[r,no head]\&\Pi_P^\tau\arrow[r]\&\{\pm1\}\arrow[r,no head]\&1\\
\&\&\tau_a\arrow[r,mapsto]\&-1
\end{tikzcd}}
(103)

mais il reste que le choix d'un foncteur (102) dépend du choix d'une orientation $\omega$ de $P$.

Un foncteur quasi-inverse de (102) s'obtient en interprétant le deuxième \unsure{mb} par la théorie de Galois des revêtements étales de $\Sigma_P^*$

\diagram{\begin{tikzcd}[ampersand replacement=\&,column sep=large]
\Pi_P\text{-ens finis}\arrow[r,"\approx"]\&
\begin{gathered}\text{revêtements ramifiés}\\
\text{finis de }\Sigma_P,\text{ étales}\\
\text{au dessus de }\Sigma_P^*
\end{gathered}
\end{tikzcd}}
(104)

si on a un tel rev. \deleted{étale} $X$ de $\Sigma_P$, l'image inverse de la carte \deleted{\unsure{qui est}} \inserted{canonique} sur $\Sigma_P$ (\unsure{jouant} le rôle de \uline{carte, $P$-pondérée orientée universelle})

\margin{93}
